%% toc: P15 172 172.7,172.8,172.9,172.10
\chapter{Sensitivity Conclusions under Changed Dependence and Formulation}
\label{sec:p15-w02b-p15-069-assumption-sensitivity}

\Cref{sec:p15-global-sensitivity} separates uncertainty propagation,
local derivatives and global sensitivity, reconstructs the argument against
one-at-a-time designs, and derives the first-order and total indices under
independent inputs. \Cref{thm:p15-global-ranking} then shows that a ranking
and a threshold can reverse when an input range changes in a fixed
independent linear model. This chapter treats the two other changes named
in \cref{def:p15-global-lookup}: a change of the joint dependence structure
with the model and the marginal laws held fixed, and a change of model
formulation with the input law and the baseline derivatives held fixed.
For each change it states which conclusions survive, which reverse, and at
which value of the assumption parameter the status changes.

Saltelli and Annoni describe a global analysis as an attempt to show that a
desired inference continues to hold while the input assumptions range over
plausible alternatives, and they cite Leamer's organized sensitivity
analysis, in which a credible neighborhood of assumptions is mapped to an
interval of inferences \cite[p.~1509, Section~2]{SaltelliAnnoni2010Sensitivity}.
They define the two indices by their Equations~(1) and~(2), recommend the
total index for fixing non-influential factors, and give computational
recipes for independent inputs only, referring the dependent case to other
work \cite[p.~1512, Section~5.4]{SaltelliAnnoni2010Sensitivity}.
The results below are derived here from the conditional-variance
definitions of \cref{def:p15-global-indices}; they are not results stated
in that article.
\index{sensitivity analysis!dependent inputs}
\index{sensitivity analysis!model formulation}
\index{factor fixing}

\section{What is being tested}

\begin{hbdefinition}{def:p15-w02b-p15-069-assumption-family}{Conclusions indexed by an assumption family}
Fix a finite index set \(I=\{1,\ldots,d\}\). An \emph{assumption family}
is a set \(\mathcal A\) whose members \(\alpha\) each specify a joint input
law \(P_\alpha\) on a Borel domain \(\mathcal X\subseteq\mathbb R^d\) and a
measurable model output \(f_\alpha:\mathcal X\to\mathbb R\), with
\(Y_\alpha=f_\alpha(X)\), \(X\sim P_\alpha\), square integrable and
\(V_\alpha=\operatorname{Var}(Y_\alpha)>0\). The indices
\(S_i(\alpha)\) and \(S_{T_i}(\alpha)\) are those of
\cref{def:p15-global-indices} computed under \(P_\alpha\) and \(f_\alpha\).

A \emph{conclusion} is a statement whose truth value is a function of
\(\alpha\). Four kinds are used below.
\begin{enumerate}
\item A first-order ranking \(S_i(\alpha)\geq S_j(\alpha)\), or its strict form.
\item A total ranking \(S_{T_i}(\alpha)\geq S_{T_j}(\alpha)\), or its strict form.
\item A threshold status: a sure bound \(f_\alpha\leq c\) on the support
 of \(P_\alpha\), or a probability bound \(P_\alpha(Y_\alpha>c)\leq q\).
\item A factor-fixing status: the claim that replacing \(X_i\) by a
 constant changes the output by a small mean-square amount.
\end{enumerate}
The \emph{status set} of a conclusion is the set of \(\alpha\) at which it
holds. A conclusion is robust over \(\mathcal A\) when its status set is
all of \(\mathcal A\). When \(\mathcal A\) is parametrized by a real
number, a \emph{breaking value} is a boundary point of the status set.
\end{hbdefinition}

\begin{hbdefinition}{def:p15-w02b-p15-069-fixing-error}{Substitution error of fixing a factor}
Suppose \(\mathcal X=\mathcal X_i\times\mathcal X_{-i}\) is a product of
Borel sets, with \(X_{-i}\) the coordinates other than \(i\), and let
\(f:\mathcal X\to\mathbb R\) be Borel measurable. For \(c\in\mathcal X_i\),
the \emph{substitution error} of fixing \(X_i\) at \(c\) is
\[
 \varphi_i(c)=E\bigl[(f(X)-f(c,X_{-i}))^2\bigr]\in[0,\infty],
\]
where \(f(c,X_{-i})\) keeps every other coordinate at its random value
under the declared joint law. This is the error of the simplified model
that a modeler would actually run after fixing \(X_i\). It is distinct
from the total-index numerator \(E[\operatorname{Var}(Y\mid X_{-i})]\),
which is an average over the conditional law of \(X_i\) given the other
inputs.
\end{hbdefinition}

\section{A dependence family with fixed marginals}

\begin{hbdefinition}{def:p15-w02b-p15-069-coupling}{The symmetric binary coupling}
For \(\rho\in[-1,1]\), let \(P_\rho\) be the law on \(\{-1,1\}^2\) with
\[
 P_\rho(X_1=x_1,X_2=x_2)=\frac{1+\rho x_1x_2}{4}.
\]
For the three-input family, let \(X_1,X_2,\varepsilon\) be independent,
with \(X_1,X_2\) fair signs on \(\{-1,1\}\) and
\(P(\varepsilon=1)=(1+\rho)/2=1-P(\varepsilon=-1)\), and put
\(X_3=\varepsilon X_1\). Write \(Q_\rho\) for the law of \((X_1,X_2,X_3)\).
\end{hbdefinition}

\begin{hblemma}{thm:p15-w02b-p15-069-coupling}{proved=annex:p15-w02b-p15-069-coupling}{Properties of the coupling}
For every \(\rho\in[-1,1]\), \(P_\rho\) is a probability law whose two
marginals are fair signs, with
\[
 E[X_2\mid X_1]=\rho X_1,\qquad
 \operatorname{Var}(X_2\mid X_1)=1-\rho^2,\qquad
 \operatorname{Corr}(X_1,X_2)=\rho,
\]
and the same identities with the roles of the coordinates exchanged.
The coordinates are independent exactly when \(\rho=0\). All four points
have positive probability exactly when \(|\rho|<1\); for \(\rho=1\) the
support is the diagonal \(x_1=x_2\), and for \(\rho=-1\) it is the
antidiagonal \(x_1=-x_2\).
Under \(Q_\rho\), all three marginals are fair signs, the pair
\((X_1,X_3)\) has law \(P_\rho\), and \(X_2\) is independent of
\((X_1,X_3)\).
\end{hblemma}

The family \(\{P_\rho\}\) changes only the dependence: each member has the
same marginal laws, and \(\rho=0\) is the independent law that a
marginal-by-marginal specification would select.
\Cref{thm:p15-global-dependent-boundary} is the member \(\rho=1\) with
output \(Y=X_1\).

\begin{hbtheorem}{thm:p15-w02b-p15-069-linear-two}{proved=annex:p15-w02b-p15-069-linear-two}{Two-factor linear models under the coupling}
Let \((X_1,X_2)\sim P_\rho\) and \(Y=a_1X_1+a_2X_2\) with real
\(a_1,a_2\), and suppose
\(V=a_1^2+a_2^2+2a_1a_2\rho>0\). Then
\[
 S_1=\frac{(a_1+a_2\rho)^2}{V},\qquad
 S_2=\frac{(a_2+a_1\rho)^2}{V},
\]
\[
 S_{T_1}=\frac{a_1^2(1-\rho^2)}{V},\qquad
 S_{T_2}=\frac{a_2^2(1-\rho^2)}{V}.
\]
\begin{enumerate}
\item Both differences satisfy
 \(V(S_1-S_2)=V(S_{T_1}-S_{T_2})=(a_1^2-a_2^2)(1-\rho^2)\).
 Hence for \(|\rho|<1\) the first-order ranking and the total ranking of
 the two factors coincide with their independent-law ranking, which
 orders \(|a_1|\) and \(|a_2|\); at \(|\rho|=1\) both rankings are ties.
\item The model is additive, yet
 \(V(S_1+S_2-1)=\rho\bigl((a_1^2+a_2^2)\rho+2a_1a_2\bigr)\), so
 \(S_1+S_2\ne1\) whenever \(\rho\ne0\) and
 \((a_1^2+a_2^2)\rho\ne-2a_1a_2\). The independent-law inequality
 \(S_i\leq S_{T_i}\) can fail: for \(a_1=1\), \(a_2=2\), \(\rho=1/2\),
 \(S_1=4/7>3/28=S_{T_1}\).
\item For \(a_1=1\), \(a_2=2\): \(P_\rho(Y>2)=(1+\rho)/4\). The probability
 bound \(P_\rho(Y>2)\leq1/10\) holds exactly when \(\rho\leq-3/5\), and the
 sure bound \(Y\leq1\) holds exactly when \(\rho=-1\).
\end{enumerate}
\end{hbtheorem}

For the model \(Y=X_1+2X_2\) of \cref{thm:p15-global-ranking}, now with
fair-sign inputs, the independent law gives
\(S_1=S_{T_1}=1/5\) and \(S_2=S_{T_2}=4/5\). At \(\rho=-3/5\) the values are
\[
 S_1=\frac1{65},\quad S_2=\frac{49}{65},\quad
 S_{T_1}=\frac{16}{65},\quad S_{T_2}=\frac{64}{65},
\]
so \(S_1+S_2=10/13\) although the model is additive.
The ranking of the two factors is unchanged throughout \(|\rho|<1\),
but the threshold conclusion \(P(Y>2)\leq1/10\) fails under independence
and holds once \(\rho\leq-3/5\). Item~1 is special to two factors with
linear conditional means; the next result shows that a third input
changes the picture.

\begin{hbtheorem}{thm:p15-w02b-p15-069-duplicate}{proved=annex:p15-w02b-p15-069-duplicate}{Dependence reverses rankings without changing the model}
Let \((X_1,X_2,X_3)\sim Q_\rho\) and \(Y=f(X)=2X_1+X_2\), so that
\(X_3\) does not enter the model. Then \(V=5\) for every \(\rho\), and
\[
 (S_1,S_2,S_3)=\Bigl(\frac45,\frac15,\frac{4\rho^2}5\Bigr),\qquad
 (S_{T_1},S_{T_2},S_{T_3})=\Bigl(\frac{4(1-\rho^2)}5,\frac15,0\Bigr).
\]
\begin{enumerate}
\item Under independence, \(\rho=0\), both rankings are
 \(X_1\succ X_2\succ X_3\), and \(S_i=S_{T_i}\) for each \(i\).
\item The first-order ranking of \(X_2\) and \(X_3\) reverses exactly when
 \(|\rho|>1/2\), with a tie at \(|\rho|=1/2\).
\item The total ranking of \(X_1\) and \(X_2\) reverses exactly when
 \(|\rho|>\sqrt3/2\), with a tie at \(|\rho|=\sqrt3/2\). At \(|\rho|=1\),
 \(S_{T_1}=0\).
\item For \(\rho\ne0\), \(S_1>S_{T_1}\).
\item \(S_{T_3}=0\) for every \(\rho\): \(Y\) is a function of
 \((X_1,X_2)\) under every member of the family.
\end{enumerate}
\end{hbtheorem}

At \(\rho=9/10\) the values are
\[
 (S_1,S_2,S_3)=\Bigl(\frac45,\frac15,\frac{81}{125}\Bigr),\qquad
 (S_{T_1},S_{T_2},S_{T_3})=\Bigl(\frac{19}{125},\frac15,0\Bigr).
\]
An inert input ranks second by first-order index, and the factor with the
largest coefficient ranks below \(X_2\) by total index. Neither reversal
reflects a change in \(f\). The first-order index of \(X_3\) measures what
\(X_3\) reveals about \(X_1\); the total index of \(X_1\) is small because,
once \(X_3\) is known, little uncertainty about \(X_1\) remains.

\section{Factor fixing under independence and under dependence}

\begin{hbtheorem}{thm:p15-w02b-p15-069-fixing}{proved=annex:p15-w02b-p15-069-fixing}{When the total index controls the cost of fixing}
Let \(\mathcal X=\mathcal X_i\times\mathcal X_{-i}\), \(f\) Borel
measurable, \(Y=f(X)\in L^2\) and \(V>0\). Write \(P_i\) for the law of
\(X_i\).
\begin{enumerate}
\item For every joint law and every Borel \(h\) on \(\mathcal X_{-i}\),
 \[
  E\bigl[(Y-h(X_{-i}))^2\bigr]\geq V\,S_{T_i},
 \]
 with equality exactly when \(h(X_{-i})=E[Y\mid X_{-i}]\) almost surely.
 In particular \(\varphi_i(c)\geq VS_{T_i}\) for every \(c\in\mathcal X_i\).
\item If \(X_i\) is independent of \(X_{-i}\), then
 \[
  \int_{\mathcal X_i}\varphi_i(c)\,P_i(dc)=2V\,S_{T_i},
 \]
 and the set of \(c\) with \(\varphi_i(c)\leq2VS_{T_i}\) has positive
 \(P_i\)-probability.
\item In the family of \cref{thm:p15-w02b-p15-069-duplicate}, with \(f\)
 defined on \(\mathbb R^3\) by the same formula, for every \(\rho\) and
 every real \(c\),
 \(\varphi_1(c)=4(1+c^2)\geq4\), while \(VS_{T_1}=4(1-\rho^2)\).
 Hence for \(\rho^2>1/2\) no constant satisfies the independent-law bound
 \(\varphi_1(c)\leq2VS_{T_1}\), and at \(|\rho|=1\) the total index is zero
 while every substitution error is at least \(4=\tfrac45V\).
 Moreover \(\tfrac12\int\varphi_1\,dP_1=4=\tfrac45V\) for every \(\rho\).
\end{enumerate}
\end{hbtheorem}

Under independence, a small total index therefore guarantees a constant
whose substitution error is at most twice the total-index numerator; item~2
is the pick-freeze identity of \cref{thm:p15-global-pick-freeze} read as a
statement about substitution. Under dependence, only the lower bound of
item~1 survives. A screening rule ``fix \(X_i\) when \(S_{T_i}<\tau\)''
applied at \(\rho=9/10\) and \(\tau=1/5\) would fix \(X_1\), whose
substitution error is at least \(4/5\) of the output variance. The last
assertion of item~3 also shows what a recipe that resamples \(X_1\) from its
marginal computes under dependence: it returns the independent-law value
\(4/5\), and it evaluates \(f\) at input combinations that have probability
zero when \(|\rho|=1\).

\section{A change of model formulation}

\begin{hbtheorem}{thm:p15-w02b-p15-069-formulation}{proved=annex:p15-w02b-p15-069-formulation}{Equal baseline derivatives, reversed global ranking}
Let \(X_1,X_2\) be independent and uniform on \([-1,1]\). For real
\(\gamma\), consider the alternative formulation
\[
 f_\gamma(x_1,x_2)=x_1+\gamma x_1^2+2x_2,\qquad (x_1,x_2)\in[-1,1]^2.
\]
Every \(f_\gamma\) has the value \(0\) and the gradient \((1,2)\) at the
baseline \((0,0)\), and \(f_0\) is the model of
\cref{thm:p15-global-ranking}.
\begin{enumerate}
\item Each \(f_\gamma\) is additive, so \(S_i=S_{T_i}\), and
 \[
  S_1=\frac{15+4\gamma^2}{75+4\gamma^2},\qquad
  S_2=\frac{60}{75+4\gamma^2}.
 \]
 The ranking \(S_2>S_1\) holds exactly when \(\gamma^2<45/4\); it is a tie at
 \(|\gamma|=3\sqrt5/2\) and reverses beyond.
\item The sure bound \(f_\gamma\leq3\) on \([-1,1]^2\) holds exactly when
 \(\gamma\leq0\). For \(\gamma>0\), \(P(f_\gamma(X)>3)>0\).
\item For \(\gamma=4\): \(S_1=79/139\), \(S_2=60/139\), the maximum of
 \(f_4\) is \(7\), and \(P(f_4(X)>3)\geq1/16\).
\end{enumerate}
\end{hbtheorem}

The two conclusions break at different places. The threshold status changes
at \(\gamma=0\), so any positive curvature in \(x_1\) breaks it; the ranking
changes only for \(|\gamma|>3\sqrt5/2\). For
\(\gamma<-3\sqrt5/2\) the ranking reverses while the threshold still holds.
Since every member has the same baseline value and gradient, no local
derivative at the baseline can discriminate between these formulations,
which is the point of the distinction drawn in
\cref{def:p15-global-model}.

\section{Procedure and a worked report}

\begin{hbmethod}{def:p15-w02b-p15-069-report}{Reporting a conclusion across alternative assumptions}
\begin{enumerate}
\item State each conclusion in one of the four forms of
 \cref{def:p15-w02b-p15-069-assumption-family}, including the threshold
 or screening level.
\item Declare the assumption family: the admissible ranges, the
 dependence structures, including the independent member when it is
 admissible, and the alternative formulations, each with its
 justification. A dependence family should keep the declared marginals
 fixed when the question is about dependence alone.
\item For each member, recompute both first-order and total indices from
 their conditional definitions under that member's joint law. Product-law
 identities, such as the decomposition of \cref{thm:p15-global-anova} and
 the pick-freeze estimators of \cref{thm:p15-global-pick-freeze}, apply
 only to independent members.
\item For factor fixing under a dependent member, report the substitution
 error \(\varphi_i\) of the simplified model that will actually be run, not
 only \(S_{T_i}\); \cref{thm:p15-w02b-p15-069-fixing} shows that the two
 can differ by an unbounded factor.
\item For each conclusion, report its status set and breaking values, and
 state separately which conclusions hold over the whole family.
\end{enumerate}
\end{hbmethod}

\begin{hbexample}{def:p15-w02b-p15-069-worked-report}{Report for the duplicated-information and curvature families}
For \(Y=2X_1+X_2\) under \(Q_\rho\), with the conclusions
(a) \(S_2\geq S_3\), (b) \(S_{T_1}\geq S_{T_2}\),
(c) ``\(X_3\) may be fixed'', and (d) ``\(X_1\) may be fixed at
screening level \(\tau=1/5\) of the total index'':
\begin{center}
\begin{tabular}{lllll}
 member & (a) & (b) & (c) & (d) by \(S_{T_1}<\tau\)\\\hline
 \(\rho=0\) & holds & holds & holds & no, \(S_{T_1}=4/5\)\\
 \(\rho=3/4\) & fails, \(\frac9{20}>\frac15\) & holds, \(\frac7{20}>\frac15\) & holds & no\\
 \(\rho=9/10\) & fails & fails, \(\frac{19}{125}<\frac15\) & holds & yes; \(\varphi_1\geq4\)
\end{tabular}
\end{center}
The status set of (a) is \(|\rho|\leq1/2\); that of (b) is
\(|\rho|\leq\sqrt3/2\); (c) holds over the whole family, since \(f\) does not
read \(x_3\). The screening rule in (d) passes exactly for
\(|\rho|>\sqrt3/2\), where \(4(1-\rho^2)/5<1/5\), but the substitution
error then remains at least \(4/5\) of \(V\), so the factor-fixing
conclusion is not supported by the total index under those members.

For \(f_\gamma\) of \cref{thm:p15-w02b-p15-069-formulation} under independent
uniform inputs, with the conclusions (e) \(S_2>S_1\) and (f)
\(f_\gamma\leq3\) throughout the square:
\begin{center}
\begin{tabular}{lcc}
 formulation & (e) & (f)\\\hline
 \(\gamma=0\) & holds, \(4/5>1/5\) & holds, maximum \(3\)\\
 \(\gamma=1\) & holds, \(60/79>19/79\) & fails, maximum \(4\)\\
 \(\gamma=4\) & fails, \(60/139<79/139\) & fails, maximum \(7\)\\
 \(\gamma=-4\) & fails & holds
\end{tabular}
\end{center}
Neither (e) nor (f) is robust over \(\gamma\in[-4,4]\). A report that gave
only the \(\gamma=0\) row would present a model-based ranking as if it did
not depend on the formulation.
\end{hbexample}

\begin{hbwarning}{def:p15-w02b-p15-069-limits}{What these families do and do not show}
The coupling and curvature families are chosen so that every index and
probability is exact. They show that reversals of the stated kinds occur,
and where; they do not estimate how often such reversals occur in applied
models. A status set is only as credible as the declared family: an
unjustified dependence parameter or formulation can manufacture a reversal,
and an over-narrow family can hide one. Under dependence the indices remain
defined, but their interpretation changes: \(S_i\) includes what \(X_i\)
reveals about correlated inputs, and the independent-law relations
\(S_i\leq S_{T_i}\) and \(\sum_iS_i=1\) for additive models fail
(\cref{thm:p15-w02b-p15-069-linear-two,thm:p15-w02b-p15-069-duplicate}).
None of these quantities identifies an intervention effect
(\cref{thm:p15-global-causal-boundary}). Estimation error for sampled
indices is a further source of apparent reversal and must be reported
separately.
\end{hbwarning}

The covariance propagation of \cref{thm:p15-measurement-affine-covariance}
answers the uncertainty question for affine models with correlated inputs;
the results here concern the attribution and decision questions that sit on
top of it. The conditional-variance identities used throughout are those of
\cref{thm:p15-conditional-variance}.
